\chapter{Tuning while training: the population-based family}
\label{sec:population}

Everything so far searches for one fixed configuration, and the second bend
of Section~\ref{sec:formal} has been waiting patiently: in the running
example, the best learning rate at step zero is not the best learning rate
at step 100 million, because the agent's own improving policy keeps changing
the data it trains on. Population-based training (PBT) is the family that
takes this seriously. It searches for a \emph{schedule}, and it does so
inside a single training run, which is why it occupies a special place in
reinforcement learning, where the optimal learning rate, entropy bonus, or
batch size at the start of training is rarely optimal at the end.

\section{PBT: exploit, explore, inherit}

PBT trains a population of $B$ agents in parallel, each with its own
hyperparameters and its own network weights \citep{jaderberg2017pbt}. Every
$t_{\mathrm{ready}}$ steps, two operations run
(Figure~\ref{fig:pbt}). In the \emph{exploit} step, each agent in the
bottom fraction of the population (by current evaluation score) abandons
its own progress: it copies both the network weights and the
hyperparameters of an agent sampled from the top fraction, a rule known as
truncation selection. In the \emph{explore} step, the copied
hyperparameters are perturbed, in the original paper by random
multiplicative jitter or resampling from the prior. Training then
continues from the copied weights. Because weights are inherited rather
than restarted, the population performs a single collective training run
whose hyperparameters adapt on the fly, and the lineage of the final best
agent traces out a hyperparameter schedule that no pointwise tuner could
have expressed: in Figure~\ref{fig:pbt}, the winning agent's history runs
through two exploit events and three different learning rates, and that
whole trajectory, not its endpoint, is the method's output.

The figure repays one slow read-through, event by event, because every
later argument about this family refers back to it. Four agents train
in parallel, each with its own learning rate. At the first ready event,
the population is ranked by current evaluated return; the bottom agent
(lane two) is culled, receives a byte-for-byte copy of the top agent's
network weights and hyperparameters, has those hyperparameters
perturbed, and resumes training from the copied weights rather than
from scratch.

Between the first and second ready events, four
\emph{different} hyperparameter settings are being evaluated while
four networks make training progress; no compute is spent on rehearsal
runs that will be thrown away, which is the family's core economy.

At
the second ready event the same surgery repeats on the new laggard.
The winning agent's bold lineage, read left to right, is a sequence of
(weights inherited, hyperparameters revised) segments: a learning-rate
\emph{schedule} discovered by selection, with no schedule ever
specified. That is the object no method in the previous chapters can
produce, and it is also, as the next paragraphs show, the source of
the family's characteristic failures.

\begin{figure}[t]
\centering
\begin{tikzpicture}[scale=1.0, font=\small]
% axes
\draw[-{Latex}] (0,0) -- (10.2,0);
\node[anchor=north east] at (10.3,-0.55) {training steps};
\node[rotate=90, anchor=south] at (-0.4,2.1) {agents (population $B=4$)};
% ready-event lines
\foreach \x in {3.0,6.0,9.0} \draw[dotted] (\x,-0.1) -- (\x,4.3);
\node[font=\footnotesize, anchor=north] at (3.0,-0.15) {$t_{\mathrm{ready}}$};
\node[font=\footnotesize, anchor=north] at (6.0,-0.15) {$2t_{\mathrm{ready}}$};
\node[font=\footnotesize, anchor=north] at (9.0,-0.15) {$3t_{\mathrm{ready}}$};
% agent lanes at y = 0.7, 1.7, 2.7, 3.7
% agent 4 (top lane): strong performer, survives; lineage bold
\draw[line width=1.5pt, blue!60!black] (0,3.7) -- (3.0,3.7);
\draw[line width=1.5pt, blue!60!black] (3.0,3.7) -- (6.0,3.7);
\draw[line width=1.5pt, blue!60!black] (6.0,3.7) -- (9.0,3.7) -- (10.0,3.7);
% agent 3
\draw[thick] (0,2.7) -- (3.0,2.7);
\draw[thick] (3.0,2.7) -- (6.0,2.7);
\draw[thick] (6.0,2.7) -- (10.0,2.7);
% agent 2: culled at t_ready, copies agent 4
\draw[thick] (0,1.7) -- (3.0,1.7);
\node at (3.0,1.7) {$\times$};
\draw[-{Latex}, thick, blue!60!black] (3.0,3.6) .. controls (3.25,2.7) .. (3.05,1.8);
\draw[line width=1.5pt, blue!60!black] (3.1,1.7) -- (6.0,1.7) -- (9.0,1.7);
\draw[line width=1.5pt, blue!60!black] (9.0,1.7) -- (10.0,1.7);
\node[blue!60!black, font=\footnotesize, anchor=west] at (3.45,2.18) {copy weights + hparams, then perturb};
% agent 1: culled at 2t_ready, copies agent 2's line
\draw[thick] (0,0.7) -- (6.0,0.7);
\node at (6.0,0.7) {$\times$};
\draw[-{Latex}, thick] (6.0,1.6) .. controls (6.25,1.1) .. (6.05,0.8);
\draw[thick] (6.1,0.7) -- (10.0,0.7);
% winner label
\node[blue!60!black, anchor=west, align=left, font=\small] at (6.45,4.15)
  {bold path: lineage of the final best agent};
\node[anchor=west, font=\footnotesize] at (0.1,4.15) {$\times$ = culled by truncation selection};
\end{tikzpicture}
\caption{One PBT run (schematic, $B = 4$). At each ready event the bottom
agents are culled and restart from copies of top agents' weights and
hyperparameters, perturbed before training resumes. The bold lineage is
the method's real product: a hyperparameter \emph{schedule} that switched
values twice mid-run, expressed by no pointwise method in this document.}
\label{fig:pbt}
\end{figure}

Two ablations in the original paper pin down where the value lives, and
they are the reason we drew the lineage rather than the final point.
Rerunning training from scratch with the final PBT-found hyperparameters
underperformed PBT itself, so the schedule rather than the endpoint
carries the gain; and populations of around twenty gave consistent
improvements with diminishing returns beyond. The cost structure is
equally distinctive: one PBT run costs $B$ training runs of wall-clock
length one, fully parallel, rather than a long sequence of trials. When a
project has $B$ machines and one calendar week, that shape is not a detail;
it is the argument.

Two failure modes matter as much as the mechanism. First, PBT is greedy:
truncation selection optimizes the \emph{current} score, and
Figure~\ref{fig:curves} already showed what current scores do to
configurations that invest now to pay off later. A small learning rate or
a heavy regularizer is exactly the blue curve of that figure, and
truncation selection is the cull arrow. FIRE-PBT (faster improvement
rate PBT) counters this by scoring
sub-populations on their rate of improvement rather than their level, and
the size of the greediness cost is worth quoting: on ImageNet, the
learning-rate schedule found by greedy PBT reached 71.7 percent top-1
accuracy where FIRE-PBT reached 76.5, essentially matching the hand-tuned
schedule at 76.6 \citep{dalibard2021firepbt}.

Second, with the modest
populations we can afford (the experiments below use $B = 8$), copying
weights across the population collapses diversity in weight space: after
a few exploit events, every agent descends from the same early winner,
and the search degenerates into local perturbation around one basin.
\citet{wan2022bgpbt} observed both effects directly: on four of their
seven Brax tasks, plain random search matched or beat both PBT and PB2 at
the same compute, which they attribute to exactly this
premature-convergence behavior.

\section{PB2: replacing random perturbation with a bandit}

The explore step is where PBT wastes its evidence: jitter treats the
population's accumulated experience as worthless. PB2 repairs it with the
machinery of Chapter~\ref{sec:bayesian} \citep{parkerholder2020pb2}. PB2
records every observed (hyperparameters, training-progress) increment
across the population and fits a Gaussian process whose kernel includes a
time coordinate, so that old evidence decays: an observation from early
training informs, but does not dictate, beliefs about what works now.

When an agent needs new hyperparameters after an exploit event, PB2
chooses them by maximizing a UCB acquisition, the rule of
\eqref{eq:ucb} applied to this time-varying surrogate: for the increment
model $\hat{f}_t$ with posterior mean $\mu_t$ and spread $s_t$, pick
\begin{equation}
\config_{\mathrm{new}}
\in \arg\max_{\config \in \cspace}\;
\mu_t(\config) + \beta_t^{1/2}\, s_t(\config),
\label{eq:pb2}
\end{equation}
maximization now because the increment is a reward. Members of the batch
whose outcomes are still pending are handled by updating the posterior
variance as if their results were already in (the fantasy device of
Section~\ref{sec:acquisition}), so the whole population can be proposed
at once.

The quantity modeled is deliberately the \emph{increment} of
performance over one interval rather than its level, which is what makes
the bandit framing clean: total performance is the sum of increments, so
maximizing final performance equals minimizing cumulative regret over
increments, and the authors prove a sublinear-regret bound for the
resulting time-varying GP bandit. Sublinear regret has a plain reading:
the average quality gap between the bandit's choices and the best
choices available at each moment goes to zero as training proceeds.

A follow-up, PB2-Mix, adds categorical hyperparameters through a
hierarchical explore step: a time-varying adversarial bandit picks the
categories, then the GP picks the continuous coordinates conditioned on
that choice, with the regret guarantee extended to the mixed space
\citep{parkerholder2021pb2mix}.

The empirical promise is sample
efficiency at small population sizes, but two independent observations
temper it: the BG-PBT ablations show a plain GP surrogate limits how much
of the space PB2 can cover \citep{wan2022bgpbt}, and in the
budget-constrained study of \citet{eimer2023hyperparameters}, PB2's tuned
configurations changed at most once over a run (so the schedule
advantage, the entire point of Figure~\ref{fig:pbt}, went unused) and its
incumbents transferred poorly to fresh seeds.

\section{BG-PBT: trust-region BO plus generational architecture search}

BG-PBT, the method we were specifically asked to evaluate, upgrades both
halves of PBT and adds architecture search \citep{wan2022bgpbt}. We read
the paper in full; what follows is its actual content rather than its
abstract. Readers will recognize every ingredient: the search space is
the running example's, the searcher is Section~\ref{sec:highdim}'s
trust-region machinery, and the schedule-versus-point distinction is
Figure~\ref{fig:pbt}'s.

The search space is the full RL configuration. For their PPO-on-Brax
experiments this means nine training hyperparameters (learning rate,
discount, entropy coefficient, unroll length, reward scaling, batch size,
update epochs, GAE $\lambda$, clipping $\epsilon$) plus six architecture
parameters (width, depth, and a spectral-normalization flag for policy
and value networks separately), the 15-dimensional mixed space of
continuous, ordinal, and categorical coordinates that
Section~\ref{sec:running} adopted as our running example.

The explore
step becomes a high-dimensional mixed-space BO agent built on
Casmopolitan \citep{wan2021casmopolitan}: kernels tailored to each
coordinate type (their extension treats ordinals as ordered rather than
categorical, with measurable gains in their ablation), acquisition
optimization that interleaves local search over discrete coordinates
with gradient steps over continuous ones, and trust regions around the
current best configuration exactly as in Figure~\ref{fig:turbo},
expanding on successes, shrinking on failures, and triggering a restart
with a fresh surrogate when they collapse, guided by a UCB rule on an
auxiliary global GP. Restarts matter more here than in the static
setting of Section~\ref{sec:highdim}: in a time-varying problem they are
what stops the surrogate from anchoring on stale evidence.

The authors
extend the regret guarantee of PB2 to this trust-region, mixed-space,
time-varying setting (their Theorem 3.1), with the usual caveat that
the guarantee covers the version without architecture transfer.

Architecture search enters through \emph{generations}, and the device
that makes it affordable is distillation
(Figure~\ref{fig:generations}). Within a generation, architectures are
fixed per agent; truncation selection copies weights \emph{and
architecture}, so poor architectures are squeezed out of the population,
and the BO agent conditions its hyperparameter suggestions on the
architecture coordinates as context. When the population's evaluated
return stagnates past a patience threshold, a new generation begins: a
fresh set of architectures is proposed (by random search plus successive
halving over the architecture subspace, later informed by a BO model fit
to the best score each architecture achieved), and the new agents are
trained initially by on-policy distillation from the best agent of the
outgoing generation, with a joint loss that anneals from imitation
(value regression plus a KL term to the teacher policy) into the plain
RL objective. Distillation is what makes changing architecture mid-run
affordable: the new network does not start from zero, it starts from an
imitation of the best agent so far.

\begin{figure}[t]
\centering
\begin{tikzpicture}[scale=1.0, font=\small]
\draw[-{Latex}] (0,0) -- (10.2,0);
\node[anchor=north east] at (10.3,-0.12) {training steps};
% generation 1 block
\draw[rounded corners, fill=blue!8] (0.2,0.4) rectangle (4.4,2.3);
\node[anchor=south west, font=\footnotesize] at (0.25,2.34) {gen.\ 1: architectures fixed};
\foreach \y in {0.7,1.2,1.7,2.05} \draw[thick] (0.5,\y) -- (4.1,\y);
% stagnation marker
\draw[dotted] (4.4,0.15) -- (4.4,2.85);
\draw[dotted] (4.4,0.15) -- (4.4,-0.35);
\node[anchor=north, font=\footnotesize] at (4.4,-0.4) {return stagnates};
% distillation arrow over the top
\draw[-{Latex}, very thick, blue!60!black] (3.9,2.45) .. controls (4.7,3.3) and (5.4,3.3) .. (6.2,2.45);
\node[blue!60!black, font=\footnotesize, anchor=south, align=center] at (5.05,3.32)
  {best agent teaches the new generation:\\imitation loss annealing into the RL objective};
% generation 2 block
\draw[rounded corners, fill=blue!8] (5.5,0.4) rectangle (9.9,2.3);
\node[anchor=south east, font=\footnotesize] at (9.85,2.34) {gen.\ 2: fresh architectures};
\foreach \y in {0.7,1.2,1.7,2.05} \draw[thick] (5.8,\y) -- (9.6,\y);
\end{tikzpicture}
\caption{BG-PBT's generational structure (schematic). Hyperparameters
adapt continuously inside a generation via the trust-region BO explore
step; architectures change only at generation boundaries, where new
networks are warm-started by distilling the outgoing best agent rather
than training from scratch. Distillation is the piece that makes
mid-run architecture change affordable.}
\label{fig:generations}
\end{figure}

The reported results, on seven Brax environments with population size 8
and 150M environment steps: BG-PBT reaches the best mean return on five
of seven tasks against tuned PPO, random search, PBT, and PB2, with the
remaining two split; against a sequential SMAC3 baseline given fifty
full runs (about six times the compute), BG-PBT still wins on four of
seven. The discovered schedules are interpretable and match folklore:
learning rates fall over training and batch sizes grow, without any
schedule having been specified. The paper's own limitations section
notes that on Humanoid the ablation without architecture search did
better, that generational transfer voids the theoretical guarantee, and
that the distillation hyperparameters are themselves fixed by hand. That
last item deserves a smile: the state of the art in hyperparameter
optimization contains, at its center, a handful of hand-tuned
hyperparameters. The regress has to stop somewhere, and knowing where it
stops is part of understanding the method.

\section{The two experiments this chapter's verdict rests on}
\label{sec:pop-evidence}

The verdict this chapter builds toward, build BG-PBT but confine the
whole population family to studies with real parallelism, rests on two
experiments. Section~\ref{sec:bo-evidence} argued that evidence should
be weighed by its design, not its abstract, so both designs are laid
out here in enough detail to be judged, and then the argument that
reconciles them is made explicitly, because at first reading they
appear to disagree.

The first experiment is the BG-PBT paper's own
\citep{wan2022bgpbt}. Seven Brax locomotion environments; PPO; a
population of eight agents; 150 million environment steps per run; the
fifteen-dimensional mixed search space of our running example, nine
training hyperparameters plus six architecture parameters. The
comparisons: PPO with tuned fixed hyperparameters, random search over
the same full space, original PBT, PB2, and, importantly, a
sequential SMAC3 baseline given fifty complete training runs, roughly
six times the compute of one population. The scoring is mean evaluated
return, and the headline results were quoted in the previous section:
best on five of seven tasks against the population-budget baselines,
still ahead on four of seven against the six-fold-compute sequential
baseline.

Three details of the design speak well of the authors.
Random search was run over the \emph{full} joint space rather than a
crippled subset, which is how the paper itself surfaced the
uncomfortable fact that random search matched or beat PBT and PB2 on
four of seven tasks. The ablations separate the contributions
(ordinal-aware kernels help; on Humanoid the version \emph{without}
architecture search did better). And the limitations section concedes
the points a referee would raise: the theory does not cover
generational transfer, and the distillation hyperparameters are set by
hand. A paper that reports where its own method loses earns more trust
than one that does not.

The second experiment is the controlled RL tuning study of
\citet{eimer2023hyperparameters}, from a different group with no
shared authors. Three algorithms (PPO, SAC, DQN) across eleven
environments, Brax tasks among them; tuners compared at matched
budgets of ten to sixty-four \emph{full training runs}, which is the
regime our routine studies actually inhabit; candidates evaluated on
several tuning seeds with final results reported on disjoint test
seeds the tuner never saw; and the tuners' own overhead measured
(minutes for most methods, about two hours for BG-PBT at the largest
budget, against days of trial compute).

The design's discipline is
the point: matched budgets stop a hungrier method from winning by
eating more, and the seed split stops a tuner from being rewarded for
overfitting luck. In this regime, DEHB was the most consistent tuner,
random search was respectable, and the population methods trailed:
BG-PBT sat behind DEHB and even random search at sixteen and
sixty-four runs, and PB2's tuned configurations changed at most once
per run, so the schedule advantage that justifies the family's
existence went unused. The authors note, fairly, that BG-PBT would
likely benefit from larger budgets.

Now the reconciliation, because a reader holding both results needs
it. The two experiments occupy different budget regimes. One
population run in the Wan design consumes eight parallel full-length
trainings, and its evolutionary machinery gets hundreds of exploit
events over 150 million steps; the Eimer design's entire budget is
ten to sixty-four full runs total, which starves a population method
of exactly the parallelism and step count its mechanism needs.

Where
the regimes overlap, the studies agree: both find random search over
the full space embarrassingly strong at modest budgets, and both find
plain PBT and PB2 fragile with small populations. So the apparent
contradiction is not one paper being wrong; it is one boundary seen
from both sides, and the boundary is the finding.

Below roughly eight
concurrent full-length workers or a few dozen full-run equivalents,
rung-based methods with a model (DEHB, or BO plus ASHA) are the
reliable choice; above it, with a schedule-shaped question, BG-PBT is
the strongest design published for that regime, on the evidence of
Brax locomotion at population sizes of eight and generous step
budgets, and its official implementation is available to build from. The package should enforce that boundary in
software rather than trust each user's optimism, and our validation
suite re-tests it on our own pinned tasks, since a boundary
established on Brax deserves a check before it governs a default.

\section{The family since 2022}

The line has stayed active, and the newer entries share one diagnosis:
PBT's remaining free parameter, the interval $t_{\mathrm{ready}}$
between evolution steps, is itself consequential and badly understood.
Evolve too often and agents are judged before their new hyperparameters
have had any effect; too rarely and the schedule cannot track the
drift it exists to track.

Multiple-frequencies PBT runs sub-populations
that evolve at different perturbation frequencies with an asymmetric
migration rule between them, reporting that vanilla PBT can fall behind
random search over long horizons, the same greediness disease FIRE-PBT
treats \citep{doulazmi2025mfpbt}. A 2026 preprint takes the opposite
tack for small populations, adjusting the interval through restarts
that reuse weights and re-seed hyperparameters by time-varying BO
\citep{chebykin2025ipbt}.

On the off-policy side (algorithms that learn from a stored
\emph{replay buffer} of past experience rather than only from the
freshest data), SEARL shares one replay
buffer across the whole population, so the meta-search rides on data the
population already paid for, a design worth remembering for our
off-policy RLlib work \citep{franke2021searl}. None of these ships in
any framework we audited; Ray provides PBT, PBT replay, and a
continuous-only PB2 (Chapter~\ref{sec:libraries}), which is exactly why
this family is where an in-house implementation adds value.

\section{What the product should take from this family}

Three design consequences follow. First, schedules are a first-class
output: our tuner must be able to return a hyperparameter trajectory,
not only a point, and our run store must record the lineage that
produced it; Figure~\ref{fig:pbt}'s bold path must be a queryable
object, not a diagram. Second, the exploit and explore steps need
checkpoint surgery (copy weights, swap optimizer state, resume), so the
execution layer must treat checkpoints as cheap, frequent, and
programmable, which is the strongest single argument for keeping Ray
Tune underneath (its PBT scheduler already exercises this machinery,
and RLlib checkpoints compose with it). Third, the population methods
and the rung methods are complements, not rivals: the honest default is
budget-dependent, and the product should encode that guidance rather
than leaving it to each user's optimism.

\section{The verdict, argued from four angles}
\label{sec:pop-angles}

The chapter's verdict, build the PB2-Mix-to-BG-PBT line as the
flagship but confine the family behind the parallelism boundary,
should survive scrutiny from four independent directions, and it does.

From mechanism: the family's failure modes were derivable before any
experiment. Truncation selection optimizes the current score, so it
must cull slow starters (the crossing curves of
Figure~\ref{fig:curves} guarantee victims exist); weight copying must
collapse diversity in small populations, since after a few exploit
events every agent descends from one early winner. Both pathologies
follow from the mechanism as stated, both were then observed
independently (FIRE-PBT's ImageNet gap for the first, the
random-search upsets in both experiments for the second), and each
upgrade in the chapter is a targeted repair of one of them. A verdict
whose supporting evidence was predicted by the mechanism is worth
more than one discovered by accident.

From evidence: two experiments with no shared authors, designed under
different incentives (one by the method's inventors, one by outside
evaluators with a budget-matched, seed-split protocol), agree
wherever their regimes overlap and disagree exactly where the budget
regimes diverge. Section~\ref{sec:pop-evidence} laid out both designs
so this claim can be checked rather than taken.

From operations: the family's cost shape, $B$ machines for one
wall-clock training run, is the shape of our actual constraint, a
cluster with idle capacity and researchers with deadlines; and the
one genuinely hard engineering requirement, live checkpoint surgery,
is machinery Ray Tune and RLlib already productionize
(Chapter~\ref{sec:architecture}). We are not building the risky part.

From decision theory: building the line is an option purchase. Our RL
work is moving toward exactly the large-parallelism,
schedule-shaped regime where the method's advantage lives; the
boundary rule caps the downside of misuse at small budgets, where the
package will steer users to DEHB instead; and declining to build
leaves schedules unreachable at any budget, which is a capability
gap, not a savings.

And the falsification criterion, stated in advance: the boundary and
the flagship are both re-tested in our validation suite, populations
against DEHB and BO-plus-ASHA on our pinned Brax task at matched
compute on both sides of the eight-worker line. If the population
line cannot beat the rung methods even in its home regime, the
flagship is demoted and this chapter gets rewritten, not defended.

\begin{decisions}
\noindent\textbf{Decisions from this chapter.} One threshold governs
everything here: can the study run eight or more full-length trainings
in parallel? Read eight as a calibrated default rather than a law: it
is our reading of where \citet{eimer2023hyperparameters}'s budgets and
BG-PBT's population regime cross on workloads that are not ours, so
the product should ship it as a configurable threshold that dispatches
to the rung methods below the line, warns rather than refuses on
override, and records the override in the study's audit trail. V08
is what would replace the default with a measured crossover on our own
tasks. Below the line the small-budget evidence favors DEHB and
BO-plus-ASHA \citep{eimer2023hyperparameters}. Above it:
\begin{itemize}[leftmargin=1.3em, itemsep=2pt]
\item \textbf{PBT} (expose, do not build). Ray Tune already ships the
reference distributed implementation; our package should surface it
with the greediness warning attached (the FIRE-PBT ImageNet numbers,
71.7 versus 76.5 percent top-1, are the one-line version of that
warning, \citealp{dalibard2021firepbt}).
\item \textbf{PB2 extended to mixed spaces} (build; moderate cost).
Ray's PB2 covers continuous knobs only; the categorical extension
(PB2-Mix's bandit over categories, or directly the mixed-space kernels
we build for the trust-region searcher) is the first genuine gap our
package fills in this family.
\item \textbf{BG-PBT} (build; the flagship; months). The strongest
published design when a real population, a schedule-shaped answer, and
architecture search are all wanted at once: best mean return on five
of seven Brax tasks against tuned PPO, random search, PBT, and PB2,
and still ahead of a sequential baseline given six times the compute
\citep{wan2022bgpbt}. The costs to state plainly: it needs the
trust-region mixed-space searcher as a prerequisite, its theory does
not cover the generational transfer that makes it shine, and below the
population threshold it loses to methods a tenth its complexity.
\item \textbf{FIRE-PBT, MF-PBT, IPBT, SEARL} (watchlist, do not build
now). Each fixes a real PBT pathology (greediness, the evolution
interval, off-policy data reuse), none ships in a framework we
audited, and none has the cross-paper evidence BG-PBT has. SEARL's
shared replay buffer is the one to revisit first when our off-policy
work scales up.
\end{itemize}
\end{decisions}

\section{What would overturn these verdicts}

The flagship's fate is tied to V08, whose demotion rule was fixed
before any code existed: if BG-PBT cannot beat DEHB and BO+ASHA on
both sides of the eight-worker boundary at matched compute, it is
demoted to an option and the chapter rewritten with the population
line's status downgraded. The eight-worker boundary itself is
empirical; if real studies consistently choose population methods
below it or consistently avoid them above it, the boundary input is
re-tuned and this section updated. PB2-Mix's dependency on the
trust-region searcher is structural; if that searcher is cut for
other reasons, PB2-Mix either finds a different explore engine or is
deferred. FIRE-PBT's hold is lifted if a study on RL tasks shows it
reliably outperforming the simpler PBT variants already included.
